\section{旧构建的实验结果}
\label{sec:archive-results}
以下结果保留原构建、存储与测量条件，属于旧版本证据；当前正文采用冻结评估构建 721800c 的三组实验。旧版性能与 CPU 测量来自不同构建和负载，不能混用为同一系统的容量推导。

\subsection{无需逐跳结算的连续花费}
\label{archive:sec:eval-continuation}

我们把认证续花与在各跳之间等待公共执行进行比较。在基线版本上，100 跳链到达最后一次认证接收的中位耗时为 158.846\,ms，而每一跳都等待公共确认时为 65.204\,s。三个钱包轮流转账 100 CAL，每笔支付都由其前一笔的输出在线构造，并用自己的最终 FUEL 支付费用。每种模式重复三次，在九服务的内存配置上共得到六次运行、600 笔支付。两种模式都以最后一次认证接收为终点，等待模式则在各跳\emph{之间}增加一次公共确认。在委员会 commit、flush 和 gossip 设置分别为 250、10 和 10\,ms 时，公共执行及其观察每跳耗时约 0.65\,s。认证续花消除了收款方路径上的这段等待。

图~\ref{archive:fig:continuation} 给出六次整链耗时。快速运行耗时 158.492--159.905\,ms。快速模式下全部 297 个后继请求，都消费其前一笔已认证的输出，并在该输出的认证接收之后、其公共完成被观察到之前发出。一个单输入、单输出的 TXCer 为 779 字节，请求在第一跳为 1,979 字节，之后为 2,766 字节。对于这种交易形态，直接证据的大小在全部 100 跳中保持不变，因为每个请求携带的是其直接输入的证书，而不是祖先链。

\begin{figure}[t]
  \centering
  \includegraphics[width=\columnwidth]{continuation-current-zh.pdf}
  \caption{基线版本上 100 跳链到最后一次认证接收的耗时。圆点为每种模式的三次独立运行；横线标出它们的中位数。对数时间轴在两种模式中包含相同的终点。只有等待式基线要求在相邻支付之间进行公共确认。}
  \label{archive:fig:continuation}
\end{figure}

为检验这一优势是否依赖于内存中的成员存储，我们在九服务的落盘配置上，使用 NoSync 钱包和用户自付 FUEL，做了三组配对的 100 跳试验。这些试验仅在成员存储上不同（NoSync 或同步；顺序为 N/S、S/N、N/S），每组之后各跟一条同步的等待链。中位耗时分别是：NoSync 续花 681.127\,ms，同步续花 4.471\,s，同步等待 65.013\,s。因此在成员提交为同步时，认证续花仍比等待快 14.5 倍。全部 900 笔支付完成，且各委员会状态一致。

在归档的交付门控版本上，认证之后有一个投递门控，它等待三名成员 INSTALL 完整材料，或等待经认证的公共执行成功。该门控不收集第二个签名法定人数；改为立即交付会缩短续花。在每秒 200 笔支付、持续 30\,s 的条件下，每种模式三次配对运行各完成 6,000 笔支付，运行级接收 P50 的中位数从有门控时的 2.167\,ms 降到 1.300\,ms（下降 40.0\%）。另外的 100 跳试验从 287.983 降到 200.356\,ms（下降 30.43\%）。授权不变，补充材料列出了配对运行。

\subsection{有资金支撑的关闭与偿还}
\label{archive:sec:eval-liabilities}

本实验检验赔付是否保持后继支付不变，以及每次赔付是否在其源交易随后执行时恰好偿还一次。

\noindent\textbf{赔付与偿还矩阵（基线版本）。}
九次运行使用随机种子 23、37 和 59，被扣留父交易的比例分别为 0\%、1\% 和 5\%。每次运行按每秒 20 个两跳单元的速率提供负载，持续 60\,s（1,200 个单元，2,400 笔支付），对应一份固定的 60,000-CAL 备付金，使用用户自付 FUEL 且不充值。九服务部署使用落盘的委员会应用状态和 BlockStore，成员和网关在内存中，钱包为 NoSync。驱动程序只有在每个被扣留父交易的赔付已在全部四个副本上提交并完成安装之后，才释放该父交易，所以每个被扣留的源交易都在 30-s 公共时间期限之后被赔付，并在其后执行。

全部 21,600 笔支付都完成关闭。216 次赔付（1\% 的每次运行 12 次，5\% 的每次运行 60 次）各自之后都有偿还。每次运行结束时，没有义务处于 Open 或 Repaired，缺口为零，净 Spent 和 Reserved 为零，备付金回到 60,000\,CAL（表~\ref{archive:tab:eval-audit}）。所采样到的最低备付金在 1\% 时为 59,800--59,900\,CAL，在 5\% 时为 59,600\,CAL。它反映的是驱动程序的释放节奏，不是备付金规模的结果，补充材料绘出了三次 5\% 运行的轨迹。正常子支付的运行级接收 P50 为 1.185--1.251\,ms。从期限起算，委员会提交（在该版本中它承载历史表示）的观察时间中位数为 1.003--1.254\,s，全部四个副本上完成修订区块体物理安装的时间中位数为 1.257--1.761\,s（P95 最高 2.015\,s）；两者都包含 250-ms 的观察轮询。每笔正常支付花费 94 FUEL，每次赔付增加 5，奖励、销毁与退款之和等于费用上限，Held 为零。

\begin{table}[t]
  \caption{赔付与偿还矩阵的收尾审计。每一列都被三个随机种子复现；每次运行有 2,400 笔支付。备付金数值以 CAL 计，费用数值以千 FUEL 计，费用上限为各笔支付最大值之和。}
  \label{archive:tab:eval-audit}
  \centering
  \small
  \begin{tabular}{@{}lrrr@{}}
    \toprule
    被扣留父交易比例 & 0\% & 1\% & 5\% \\
    \midrule
    赔付 / 偿还次数 & 0 & 12 & 60 \\
    赔付 / 偿还 CAL & 0 & 1,200 & 6,000 \\
    最终备付金 & 60,000 & 60,000 & 60,000 \\
    未关闭 / 净 Spent & 0 & 0 & 0 \\
    费用上限总额 & 2,400 & 2,400 & 2,400 \\
    Rewards & 201.600 & 201.660 & 201.900 \\
    Burned & 24.000 & 24.000 & 24.000 \\
    Refunded & 2,174.400 & 2,174.340 & 2,174.100 \\
    Remaining Held & 0 & 0 & 0 \\
    \bottomrule
  \end{tabular}
\end{table}

账本层面的回归测试直接覆盖到达顺序：四跳链的全部 24 种到达顺序，在有到期赔付与无到期赔付两种情形下，各取三种签发组织序列，共得到 144 条轨迹；另有 48 条拆分--合并轨迹，在两种赔付模式下，改变一个 40/60 拆分、两条分支及其合并的全部到达顺序。192 条轨迹中的每一步都保持独立计算的资产投影守恒。每条轨迹结束时，各签发组织的备付金恢复，Spent 与 Reserved 为零，$\mathsf{Paid}=\mathsf{Rec}$，已消费的输入仍保持已消费。

两个十四服务的对照实验区分永久损失与延迟偿还。父交易始终不到达时，孙支付保留其 100 CAL，而负责的签发组织的备付金仍然低 100 CAL。当跨组织链以孙支付最先到达的顺序到达时，中间源交易关闭其自己的直接义务，根签发组织被赔付的 100 CAL 在根交易其后执行时得到偿还。两种情况下四个委员会状态一致，缺席的父交易自身的费用和工作预留可能保持未清偿。

三组配对的十四服务运行使用基线版本的初始构建（在批处理结果标签分离之前），每次 4,000 笔支付、速率 100 笔/秒，其正常运行与扣留源交易运行最终的所有者持有和备付金余额完全一致，与主文的 CAL 延迟中性定理 一致。并发赔付使接收 P95 升高为 1.98--2.21 倍，这些配对不能区分造成这一升高的具体环节。

一个使用 NoSync 钱包的存储对照把成员提交单独隔离出来：NoSync 成员时接收 P50 为 1.0\,ms，同步提交时为 43--52\,ms。

早期版本的资本系列（每个条件一次运行、没有赔付）表明，10,800--19,700\,CAL 的外部新增，足以维持从 3,600-CAL 起步、各含 6,000 个两跳单元的 300-s 运行，而固定 3,600-CAL 的对照只关闭了 6,000 个单元中的 46 个。

\subsection{跨注册组织的续花}
\label{archive:sec:eval-network}

为测量通信距离的影响，我们在补充材料所记录的一个归档版本上运行了双组织套件：两个钱包进程各持 32 个身份，经 64 个 lane 发送，每个组织提供 50 笔/秒，24 次运行共完成 342,000 笔支付。延迟比较使用十跳的同组织与跨组织负载，注入的 HTTP RTT 为 0、20 和 100\,ms，每个条件重复三次。每个代理方向增加一半的 RTT，组织内部认证和委员会 P2P 不注入延迟。

\begin{figure}[t]
  \centering
  \includegraphics[width=\columnwidth]{network.pdf}
  \caption{ACK 延迟（ms）与注入的跨站点 HTTP RTT（ms）的关系。每个点是前 120-s 样本组的三个运行级 P50 的中位数；误差条跨越这些 P50 的最小值与最大值。总提供速率为 100 笔/秒。委员会 P2P 未加延迟。}
  \label{archive:fig:network}
\end{figure}

对共同的前 120-s 样本组，三个运行级跨组 ACK P50 的中位数依次为 1.632、21.658 和 101.967\,ms，同组数值范围为 1.615--1.744\,ms（图~\ref{archive:fig:network}）。跨组 ACK P50 随注入 RTT 变化，测得的超出部分为 1.6--2.0\,ms。ACK 包含返回路径，但收款方可以在接受之后即刻发起其后继交易，无需等待该 ACK 到达前一个发送方。在 100-ms 跨组 RTT 下，三次 300-s 运行各完成 30,000 笔支付；待处理峰值为 106--110，并能排空，没有持续的尾部增长。该实验在一台物理主机内隔离出 HTTP 路径延迟。

\subsection{运行容量与服务可用性}
\label{archive:sec:eval-operating}

\noindent\textbf{独立输入吞吐（基线版本）。}
三次各 30,000 笔支付的运行以 2,000 笔/秒为目标，使用九服务的内存配置、NoSync 钱包和经授权的组织代付 FUEL。成员进程使用 200\% 的 Go GC 目标和 16 个运行时处理器。全部支付完成，四个委员会状态一致。含排空的完成速率为 1,750.27、1,820.91 和 1,887.31 笔/秒，历时 15.9--17.1\,s，接收 P50 为 2.646、2.842 和 2.824\,ms，P95 为 47.9、49.1 和 48.0\,ms。驱动程序记录到有 6,551--7,570 笔支付需要等待快速接收准入名额，P95 派发滞后为 179--229\,ms，所以提供的负载并不是均匀的 2,000 笔/秒。这些是基线版本上的短时完成测量，没有在 decision-first 版本上重新测量；该版本在普通支付上的开销由上文的六次运行对比检验。补充材料说明了 150,000 笔支付系列所用的版本和时长。

在归档版本上，冗余认证与已 INSTALL 的证据在所述故障下维持服务。九次每秒 200 笔的运行完成 216,000 笔支付，其间有一个成员被延迟或暂停 30\,s。网关丢失对照中，一名成员持有经 INSTALL 复制的完整材料，备用提交在 INSTALL 之后 2.503--2.508\,s 出现。冲突对照既没有出现相互冲突的双 QC，也没有重复的本金或费用效果，而部分批准仍然是一种额度代价。

补充材料第~S2 节的多身份测量报告每笔支付的 CPU 和 HTTP 负载开销；2,048 个身份共享客户端。历史内存增长与积压分别报告。该节还给出 LND regtest 参照，其完成事件、持久化设置和负载并发度与 Next 的测量不同。
